\documentclass[11pt,a4paper]{article}
\usepackage[margin=2.5cm]{geometry}
\usepackage{booktabs,hyperref,graphicx,array,tabularx,microtype}
\usepackage[T1]{fontenc}
\usepackage[utf8]{inputenc}
\newcommand{\sidebar}[1]{\par\smallskip\noindent\textbf{#1}\ }
\title{Warden: A Validation-Gated, Self-Improving Security Gateway\\for AI Agents and MCP Tools}
\author{Sunny Aaron}
\date{\today}
\begin{document}
\maketitle

\begin{abstract}
Warden inspects the three channels an attacker can reach in the Model Context Protocol (MCP):
tool descriptions, arguments and outputs. It tracks session taint and enforces a signed,
hash-bound policy. A red/blue loop proposes defences from confirmed bypasses, and an eight-stage
validation gate checks that each one is harmless to benign work before it may enforce. We report
measured results with confidence intervals, including negative ones. \textbf{The end-to-end
attack-success numbers (RQ1--RQ3) use a simulated victim agent, not a real language model;} read
them as properties of the defence pipeline, not as estimates of real-world risk
(Section~\ref{sec:limits}). On the locked test split, scored once, the cascade with a from-scratch
classifier blocks 98.2\% of in-distribution attacks (Wilson 95\% CI 90.5--99.7) and 88.9\% of
out-of-distribution attacks (67.2--96.9) at a 1.9\% benign block rate. A fine-tuned DeBERTa
classifier \emph{diverged during training} in the locked run, so its locked row is a broken model;
a corrected configuration reaches F1~0.98 on the development hold-out only.
\end{abstract}

\section{Introduction}
An agent that uses tools over MCP trusts three things it should not: the descriptions of the tools
it is offered, the arguments a possibly manipulated model decides to send, and the outputs the tools
return. A malicious server can poison a description; a hijacked model can exfiltrate data or take
harmful external actions; and tool outputs carry outside content (emails, web pages, tickets) that
may hold indirect prompt injections.

Warden is a gateway that sits between the agent and its MCP servers and inspects all three
channels on every message. Its defining claim is about how it \emph{changes}: a self-improving
defence is only safe if each new rule is proven not to harm benign work before it can enforce. The
contributions are:
\begin{itemize}
  \item a typed policy IR with a trusted compiler and Ed25519-signed bundles whose approvals are bound to the hash of a rule \emph{including its mode};
  \item a gateway with description pinning, spotlighting and set-only session taint;
  \item a four-layer detection cascade (L0--L3);
  \item an eight-stage validation gate (G1--G8) that a red/blue learning loop must pass; and
  \item an evaluation harness with Wilson intervals, exact McNemar tests, and a write-once locked test split.
\end{itemize}

\section{Threat model}
The protected assets are the organisation's data, its ability to take correct external actions
(email, payments), and the agent's integrity of purpose. The agent's model is \emph{not} trusted to
resist manipulation; the gateway and its signed policy bundle form the trusted computing base.
Tool outputs are untrusted data on every channel, and servers marked untrusted taint the session
on any output. We consider four adversaries:
\begin{enumerate}
  \item a malicious or compromised MCP server (poisoned or shadowing descriptions, hidden-Unicode instructions, rug pulls);
  \item an attacker who can place text where the agent will read it (a ticket, an email, a web page);
  \item a hijacked model attempting exfiltration or harmful actions with plausible-looking arguments; and
  \item an adversary against the learning loop itself, who crafts attacks whose obvious defence would also fire on benign traffic, trying to induce self-denial-of-service (RQ2).
\end{enumerate}
Out of scope are compromise of the host, signing key or trusted compiler; a malicious operator
approving a harmful rule at G8; the model provider's own safety; and network-level attacks on TLS.
All red-team activity targets only Warden's own lab, and all data in it is synthetic. The full
model is in \texttt{docs/THREAT\_MODEL.md}, with an OWASP ASI/MCP Top~10 mapping in
\texttt{docs/OWASP\_MAPPING.md}.

\section{Design}
\sidebar{Policy IR and compiler.} Rules have four modes: \texttt{off}, \texttt{shadow},
\texttt{canary} and \texttt{enforce}. The trusted compiler performs the G1 static checks, rejecting
unsafe regular expressions and malformed matchers. Bundles are Ed25519-signed, and an approval
references the canonical-JSON SHA-256 of the rule with its mode, so a shadow-mode approval cannot be
replayed to authorise enforcement.

\sidebar{Gateway.} The gateway proxies the real servers, pins tool descriptions and re-lists
upstream tools to catch silent rug pulls, applies spotlighting, and tracks session taint. Taint is
set-only: once outside content has been read, the session stays tainted, so later high-risk actions
are gated regardless of intervening clean steps. High-risk channels fail closed when a detector is
unavailable; read-only channels fail open.

\sidebar{Detection cascade.} L0 is a deterministic ruleset. L1 is embedding $k$NN over the
development corpus. L2 is a classifier: a logistic-regression model that is always available, or an
ONNX model (a from-scratch network or a fine-tuned transformer). L3 is a local LLM judge consulted
only in the uncertainty band. Short-text abstention and a contrastive benign corpus keep the
false-positive rate down.

\sidebar{Validation gate.} A proposed rule must pass, in order: G1 static checks; G2 replay over
recorded benign tool traffic (it must not fire on legitimate messages); G3 a utility budget; G4
efficacy against the confirmed bypass that motivated it; G5 generalisation to held-out variants; G6
a latency budget; G7 rule/LLM/token rate budgets; and G8 a human signature on the enforce-mode
bundle. Machine approval exists only as an explicit, audited opt-out used by CI and the lab.

\sidebar{Learning loop.} A red team mutates attacks against the live gateway in the lab. A
confirmed bypass seeds a blue team that proposes a rule, and the gate decides whether it may
promote. Defence memory records what each rule answers.

\section{Evaluation}
\subsection{Setup}
Baselines are B0 (direct, no Warden), B1 (spotlighting only), B2 (classifier on outputs only) and
B3 (the full baseline ruleset); B4 adds the rules the loop promoted. Proportions carry Wilson 95\%
intervals; paired detector comparisons use the exact McNemar test; overhead differences use a
bootstrap. The corpus is synthetic and written for this project. It has a development split, used
for all tuning and training, and two locked splits, \texttt{test\_iid} and \texttt{test\_ood}, held
under a SHA-256 lock that refuses to load without an explicit unlock flag (ADR-0009).

\textbf{Simulated victim.} RQ1--RQ3 use a scripted victim agent in place of a language model. It
follows each benign task's plan and, when it sees an attacker's instruction intact in a tool output
or description, complies with a probability taken from the scenario file (scaled down inside
Warden's spotlight boundary). The pipeline (gateway, detection, policy, lab state) is exercised end
to end while the model's susceptibility is a parameter, not a measurement. Attack-success rates therefore compare defences against each other under the same
victim; they do not estimate how often a real model would be hijacked.

\subsection{RQ1: detection and benign utility (simulated victim, development set)}
Table~\ref{tab:rq1} reports attack success rate (ASR) over $n=1159$ attack attempts and the
benign false-positive rate over 62 benign messages. The full baseline B3 cuts ASR from 31.9\% to
3.0\% (95\% CI 2.2--4.2) with 96.7\% utility under attack, which meets the target we set as hypothesis H1
(ASR $\le 5\%$, utility loss $\le 2$ points on benign tasks). Utility under attack is lowest for B2
(80.0\%); we have not isolated why.
Every baseline differs from B0 in ASR at $p<10^{-50}$ (exact McNemar), with no case where a
baseline fails an attack that B0 blocks. These are development-set numbers: the rules were written
while looking at this corpus, so they are optimistic.
\begin{table}[h]
\centering\small
\caption{RQ1 on the development set with the \emph{simulated} victim. Wilson 95\% intervals.}\label{tab:rq1}
\begin{tabular}{lll}\toprule Baseline & ASR & Wilson 95\% CI \\\midrule
B0 & 31.9\% & [29.3, 34.7] \\
B1 & 17.4\% & [15.4, 19.7] \\
B2 & 16.6\% & [14.5, 18.8] \\
B3 & 4.5\% & [3.4, 5.8] \\
\bottomrule\end{tabular}

\end{table}

\subsection{Locked test split, scored once}
\IfFileExists{tables/locked.tex}{%
The locked splits were scored exactly once, by a CI job, after the classifiers had been trained on
the development split only. The scorer refuses to overwrite its output file, so a second look is not
possible without deleting a committed result. Table~\ref{tab:locked} reports the cascade with
three different L2 layers, and each L2 alone. \texttt{test\_iid} has 56 attacks and 53 benign
messages; \texttt{test\_ood} has 18 attacks in two held-out families (compliance and encoded
injections). ``Blocked'' means the detector would stop the message; ``flagged'' includes
lower-confidence alerts (any verdict other than benign).
\begin{table}[h]
\centering\scriptsize
\caption{Locked test split, scored once. Hits/$n$ = rate [Wilson 95\% CI]. No parameter was changed after this run.}\label{tab:locked}
\begin{tabular}{lllll}\toprule
Variant & iid attacks blocked & iid benign blocked & iid benign flagged & OOD attacks blocked \\\midrule
Cascade (lite L2) & 44/56 = 78.6\% [66.2, 87.3] & 1/53 = 1.9\% [0.3, 9.9] & 6/53 = 11.3\% [5.3, 22.6] & 7/18 = 38.9\% [20.3, 61.4] \\
Lite L2 alone & 16/56 = 28.6\% [18.4, 41.5] & 0/53 = 0.0\% [0.0, 6.8] & 12/53 = 22.6\% [13.5, 35.5] & 0/18 = 0.0\% [0.0, 17.6] \\
Cascade + DeBERTa L2 & 42/56 = 75.0\% [62.3, 84.5] & 1/53 = 1.9\% [0.3, 9.9] & 6/53 = 11.3\% [5.3, 22.6] & 7/18 = 38.9\% [20.3, 61.4] \\
DeBERTa L2 alone & 0/56 = 0.0\% [0.0, 6.4] & 0/53 = 0.0\% [0.0, 6.8] & 53/53 = 100.0\% [93.2, 100.0] & 0/18 = 0.0\% [0.0, 17.6] \\
Cascade + scratch L2 & 55/56 = 98.2\% [90.5, 99.7] & 1/53 = 1.9\% [0.3, 9.9] & 1/53 = 1.9\% [0.3, 9.9] & 16/18 = 88.9\% [67.2, 96.9] \\
Scratch L2 alone & 53/56 = 94.6\% [85.4, 98.2] & 1/53 = 1.9\% [0.3, 9.9] & 2/53 = 3.8\% [1.0, 12.8] & 16/18 = 88.9\% [67.2, 96.9] \\
\bottomrule\end{tabular}

\end{table}

\sidebar{The from-scratch L2 works.} With the scratch classifier the cascade blocks 98.2\% of
in-distribution attacks and 88.9\% of out-of-distribution attacks at a 1.9\% benign block rate and
a 1.9\% benign flag rate. Cascade and classifier alone are nearly equal, so on this corpus the
classifier is doing most of the work. The intervals are wide: with 18 OOD attacks the lower bound is
67.2\%, and the corpus is small and synthetic, so these numbers say the approach is viable, not that
it will hold on real traffic.

\sidebar{The DeBERTa row is a broken model, not a measurement of DeBERTa.} Fine-tuning
\texttt{microsoft/deberta-v3-xsmall} on the 222-example training slice for three epochs gave a
validation accuracy of 55.4\% with precision, recall and F1 all exactly 0.0: it predicted the benign
class for every held-out development input. On the locked split its ONNX export flagged all 53
benign messages and blocked none of the attacks. We first suspected under-training or an export bug.
Later diagnostic runs showed otherwise: under the lockfile's PyTorch~2.15 and transformers~5.21 the
training diverges to NaN (in four diagnostic runs the loss or gradient norm was non-finite within the
first five steps, and anomaly detection reported \texttt{LogSoftmaxBackward0} returning NaN on the first
backward pass; steps also took over a minute each). A NaN model predicts class~0 for
every input, which matches the development result, and NaN scores at serving time read as
``not benign'', which matches the locked result. We did not isolate whether PyTorch or transformers
is responsible. With PyTorch~2.8.0 and transformers~4.56.2, now pinned in the \texttt{train} extra,
the same recipe (six epochs, learning rate $5\times10^{-5}$, 96 tokens) trains normally: on the
56-example development hold-out it reaches accuracy 98.2\%, precision 100\%, recall 96\% and F1~0.98,
and the served ONNX bundle matches PyTorch to $5\times10^{-7}$. \textbf{That model has not been
scored on the locked split.} The split was spent on the first run, and scoring the corrected model
on it would invalidate the split; a locked-split result for DeBERTa needs newly collected test data.
The locked row above therefore stands as the result for a model that had diverged.

\sidebar{Not yet established.} The L3 LLM judge was not run for the locked split (\texttt{judge:
false}) and no external embedding model was configured, so the table measures the cascade without L3.
}{%
The locked splits have not been scored; this section is a placeholder until the verification
workflow commits \texttt{docs/results/locked-test.json}.}

\subsection{RQ2: does the gate prevent self-denial-of-service?}
We ran the loop for five rounds of 200 attempts with half of the attacks blending in benign
phrasing, once gated and once with every proposal promoted without a gate (the RQ2 adversary).
Table~\ref{tab:rq2} summarises the result. The gate rejected 20 proposals and promoted 5. The
ungated loop promoted 20, one of which fails the benign-replay stage (it would fire on
legitimate messages). The mechanism therefore works as
designed: the G2 benign-replay stage caught the over-blocking rule that the ungated loop let
through.

\textbf{However, the damage was not measurable.} Benign utility did not move for either arm (13 of 13
tasks both before and after), so the 10$\times$-reduction hypothesis (H2) is met only by its
fallback criterion (``the ungated loop promoted a harmful rule and the gated one did not''), not by a
ratio of utility lost. The benign task set has only 13 tasks, so it cannot resolve small losses, and
the simulated victim does not react to a benign request being blocked the way a real agent would.
This is one seed; we make no claim about the size of the induced self-DoS.
\begin{table}[h]
\centering\small
\caption{RQ2: gated versus ungated loop (5 rounds $\times$ 200 attempts, blend rate 0.5, seed 7, simulated victim).}\label{tab:rq2}
\begin{tabular}{lrr}\toprule
& Gated & Ungated \\\midrule
Rules promoted & 5 & 20 \\
Gate rejections & 20 & -- \\
Harmful promotions & 0 & 1 \\
Benign utility lost (points) & 0.0 & 0.0 \\
\bottomrule\end{tabular}

\end{table}

\subsection{RQ3: ASR across loop rounds}
Table~\ref{tab:rq3} shows the gated loop over five rounds of 200 attempts each (no blending, seed 7,
simulated victim). The bypass rate peaks at 6.0\% in round 1 (12 of 200), when the attacker's
mutations first find gaps, then falls to 4.0\%, 0.5\% and 0.0\% as rules promote. Of 22 proposals,
5 were promoted and 17 were rejected by the gate, and benign utility stayed at 13 of 13 throughout.
The fall is an in-sample result: the attacker mutates from the same seed attacks the defender has
just been shown, and the later rounds have few bypasses, so the final 0.0\% has a wide interval
(Wilson upper bound about 1.9\% at $n=200$). A single seed also cannot separate the loop's effect
from the attacker running out of new mutations.
\begin{table}[h]
\centering\small
\caption{RQ3: gated loop ASR per round (simulated victim, seed 7).}\label{tab:rq3}
\begin{tabular}{rrrrrr}\toprule
Round & Attempts & Bypasses & ASR & Promoted & Gate rejections \\\midrule
0 & 200 & 1 & 0.5\% & 1 & 0 \\
1 & 200 & 12 & 6.0\% & 4 & 8 \\
2 & 200 & 8 & 4.0\% & 0 & 8 \\
3 & 200 & 1 & 0.5\% & 0 & 1 \\
4 & 200 & 0 & 0.0\% & 0 & 0 \\
\bottomrule\end{tabular}

\end{table}

\section{Limitations}\label{sec:limits}
\begin{itemize}
  \item \textbf{Simulated victim.} RQ1--RQ3 do not use a real language model. They show the pipeline behaves as designed; they do not estimate real attack success.
  \item \textbf{Small synthetic corpus} written for this project. The test splits hold 109 in-distribution and 18 out-of-distribution messages, so intervals are wide, and synthetic text is easier than real traffic.
  \item \textbf{One locked run, several variants.} Six detector variants were scored on the same locked data in a single pass. They are not corrected for multiple comparison, and the locked data cannot be reused.
  \item \textbf{No locked-split result for a working DeBERTa.} The locked DeBERTa row is a diverged model; the corrected one is evaluated on a 56-example development hold-out only.
  \item \textbf{RQ2 magnitude is unmeasured} (13-task benign set, one seed).
  \item \textbf{Partial OWASP coverage}, with gaps listed in \texttt{docs/OWASP\_MAPPING.md}; no evaluation on public benchmark runs.
  \item \textbf{Scope.} Warden sees what reaches an MCP tool call; it cannot see reasoning that never produces one.
  \item \textbf{Deployment.} The Helm chart installs on a \texttt{kind} cluster in CI and Terraform passes \texttt{validate}; neither was applied to a real cloud account, and there is no hosted deployment.
\end{itemize}

\section{Related work}
Warden builds on, and is complementary to, benchmark and defence work in this area:
AgentDojo~\cite{debenedetti2024agentdojo} for evaluating injection attacks on tool-using agents; Llama Guard and LlamaFirewall~\cite{chennabasappa2025llamafirewall} and Prompt Guard for
classifier-based screening; spotlighting~\cite{hines2024spotlighting}; and the OWASP Top~10 for
agentic applications and for MCP. The new element here is not the detectors but the validation gate
for the defence's own updates, which is the part the self-DoS experiment targets.

\begin{thebibliography}{9}
\bibitem{debenedetti2024agentdojo} E.~Debenedetti et al. AgentDojo: A dynamic environment to evaluate prompt injection attacks and defenses for LLM agents. NeurIPS Datasets and Benchmarks, 2024.
\bibitem{chennabasappa2025llamafirewall} S.~Chennabasappa et al. LlamaFirewall: An open source guardrail system for building secure AI agents. 2025.
\bibitem{hines2024spotlighting} K.~Hines et al. Defending against indirect prompt injection attacks with spotlighting. 2024.
\end{thebibliography}

\end{document}
